% \clearpage
% ////////////////////////////////////////////////////////////////////////////////////////////
\section{The individual infectiousness profile}
% \subsection{Expressing the Hudson--Kermack--McKendrick model in terms of individual infectiousness}


% \begin{figure}[ht]
%     \centering
%     % \includegraphics[width=.5\textwidth]{figures/stock.png}
%     \newcommand{\panelbox}[2]{\framebox{\begin{minipage}[c][3.2cm][c]{0.3\textwidth}\centering\textbf{#1}\\\smallskip\small #2\end{minipage}}}%
%     \begin{tabular}{@{}ccc@{}}
%     \panelbox{A}{Smooth profile\\schematic} &
%     \panelbox{B}{Stepwise profile\\schematic} &
%     \panelbox{C}{Spike profile\\schematic} \\[4pt]
%     \panelbox{D}{Cumulative infections\\(smooth)} &
%     \panelbox{E}{Cumulative infections\\(stepwise)} &
%     \panelbox{F}{Cumulative infections\\(spike)} \\[4pt]
%     \panelbox{G}{Daily new infections\\(smooth)} &
%     \panelbox{H}{Daily new infections\\(stepwise)} &
%     \panelbox{I}{Daily new infections\\(spike)}
%     \end{tabular}
%     \caption{{\bf The individual infectiousness profile impacts stochastic epidemic trajectories.} (A-C) Schematic diagrams of three individual infectiousness profile families that yield equivalent SEIR-type mean-field dynamics: the smooth (A), stepwise (B), and spike (C) profiles (Eqs.~\ref{eq:a_seir_stepwise}, \ref{eq:a_seir_smooth}, \ref{eq:a_seir_spike}). (D-I) Deterministic epidemic trajectories (blue) are numerical solutions of the mean-field renewal equation model (Eq.~\ref{eq:foi}), equivalent to a standard SEIR model with latent period $1/\gamma_{\text{SEIR}} = XX$, recovery period $1/\alpha_{\text{SEIR}} = XX$, and reproduction number $R_0 = 6$. For each individual infectiousness profile family, 1,000 stochastic simulations were generated (grey), with the cumulative infections (D-F) and daily new infection counts (G-I) depicted.}
%     \label{fig:seir_sims}
% \end{figure}

%////////////////////////////////////////////////////////////////////////////////////////////

\subsection{Components of the individual infectiousness profile}

The individual infectiousness profile $a_i(\tau)$ can be decomposed into three components that separate intensity, biological timing, and contact-related timing:
%
\begin{equation}
    a_i(t_i, \tau) = \nu_i \, b_i(\tau) \, c_i(t_i + \tau)
    \label{eq:abc}
\end{equation}
%
Here, $\nu_i$ is person $i$'s individual reproduction number (the expected number of secondary infections under a time-uniform reference contact process); $b_i(\tau)$ is a probability density capturing the biological timing of person $i$'s infectiousness under that reference process; $c_i(t)$ is a nonnegative contact modulator describing person $i$'s proportional deviation from the reference contact level at calendar time $t$; and $t_i$ is person $i$'s infection time. The contact modulator admits a useful interpretation: $\nu_i c_i(t)$ is the expected number of secondary infections person $i$ would produce if all their infectiousness were concentrated at time $t$ (\ie, if $b_i$ were a delta function at $t$). For $\nu_i$ to retain its meaning, we require $E_t[c_i(t)] = 1$, which is satisfied by any periodic function with mean 1 or any stationary random process with mean 1. Setting $c_i(t) \equiv 1$ recovers the standard decomposition $a_i(\tau) = \nu_i \, \xi_i(\tau)$ of the individual infectiousness profile into an individual reproduction number and an individual intrinsic generation interval distribution, with $b_i(\tau) = \xi_i(\tau)$.

This individual-level decomposition extends to the population level. Assuming that $c_i$ is independent of $\nu_i$ and $b_i$ (while allowing $\nu_i$ and $b_i$ to depend on one another), the population-level infectiousness profile is
%
\begin{equation}
    A(t, \tau) = E_i[a_i(t,\tau)] = E_i[\nu_i \, b_i(\tau)] \cdot E_i[c_i(t + \tau)]
    \label{eq:A_def_general}
\end{equation}
%
Defining the population-mean contact modulator $C(t) = E_i[c_i(t)]$ and the $\nu_i$-weighted average biological profile $B(\tau) = E_i[b_i(\tau) \, \nu_i / R_0]$, this becomes
%
\begin{equation}
    A(t, \tau) = R_0 \, B(\tau) \, C(t + \tau)
    \label{eq:ABC}
\end{equation}
%
Since each $c_i$ has time-average 1, the population-mean $C(t)$ also has time-average 1, so $C$ does not inflate or deflate the overall level of transmission. When $C(t) \equiv 1$ (\ie, contacts are time-uniform), we recover $A(\tau) = R_0 \, B(\tau) = R_0 \, g(\tau)$, where $B(\tau)$ is the intrinsic generation interval distribution. In this paper, we fix $\nu_i = R_0$ for all individuals, isolating the effect of biological-timing variation from the previously studied effects of variation in the individual reproduction number \citep{lloyd-smith_superspreading_2005}. Under this constraint, $B(\tau) = E_i[b_i(\tau)] = g(\tau)$.

% ////////////////////////////////////////////////////////////////////////////////////////////

\subsection{A time-shifted Gamma model for biological infectiousness}

Our central goal is to examine how the shape of the individual infectiousness profile $a_i(\tau)$ impacts epidemic dynamics, holding all else (including $R_0$ and $g(\tau)$) constant. The ``stepwise'' expression for $a_i(\tau)$ in Eq.~\ref{eq:a_seir_stepwise} poses a difficulty: it introduces inherent overdispersion in the secondary infection distribution ($\nu_i = \int_0^\infty a_i(\tau) d\tau \sim XX$), while for the smooth and spike profiles, $\nu_i$ is uniform across individuals ($\nu_i = R_0 \forall i$). This overdispersion affects stochastic epidemic dynamics, including the probability of epidemic establishment \citep{lloyd-smith_superspreading_2005}. The impact of overdispersion on stochastic epidemic dynamics has been studied extensively elsewhere \citep{lloyd-smith_superspreading_2005}. Here, we seek a family of functions that do not feature this inherent overdispersion, using the smooth and spike infectiousness profiles as book-ends.

Let us constrain the generation interval distribution $g(\tau)$ to be unimodal, and let us seek a family of functions $b_i(\tau)$ that are also unimodal with the narrowness of the function (variance of the distribution) modulated by a parameter $\psi \in [0,1]$, where $\psi \rightarrow 0$ yields a delta function (the spike case) and $\psi \rightarrow 1$ makes $b_i(\tau) \equiv g(\tau) \forall i$ (the smooth case). If we assume $g(\tau)$ is Gamma-distributed, which is a common choice for modeling generation interval distributions, the setup is straightforward. Specifically, let $g(\tau)$ be the $\text{Gamma}(\alpha, \beta)$ pdf, where $\beta = \alpha / \bar{g}$. Here, $\bar{g}$ is the mean generation time and $\alpha$ controls the shape of the generation interval distribution. A key property of the Gamma distribution is that the sum of two Gamma-distributed random variables with shape parameters $\alpha_1$ and $\alpha_2$ and common rate $\beta$ is also Gamma-distributed with shape $\alpha_1 + \alpha_2$ and rate $\beta$. Let us assume now that the individual-level biological infectiousness profile $b_i(\tau)$ follows a shifted Gamma distribution preceded by a latent period $l_i$, which is also Gamma-distributed, \ie:
%
\begin{equation}
    b_i(\tau) = \begin{cases}
      0 & \tau \leq l_i  \\
      f_\epsilon(\tau - l_i) & \tau > l_i
    \end{cases}
    \label{eq:gamma_timeshift}
\end{equation}
%
where $l_i \sim \text{Gamma}((1-\psi)\alpha, \beta)$ and $f_\epsilon(\tau)$ is a $\text{Gamma}(\psi \alpha, \beta)$ pdf. The parameter $\psi \in [0,1]$ determines the punctuatedness of $b_i(\tau)$, with $\psi \rightarrow 0$ making $b_i(\tau)$ a delta function concentrated at $\tau = l_i$, and $\psi \rightarrow 1$ making $b_i(\tau) = g(\tau)$. As such, this provides a continuum between the spike and smooth cases described in the Introduction. Due to the additive property of Gamma distributions with equal rate, these individual-level profiles are guaranteed to converge to the population-level $g(\tau) \sim \text{Gamma}(\alpha,\beta)$. Furthermore, this convergence does not depend on the entire population sharing the same $\psi$-value; indeed, unique $\psi_i$ may be chosen from any distribution with support on $[0,1]$ and the convergence holds ({\bf Supplementary Methods}).

% ////////////////////////////////////////////////////////////////////////////////////////////

\subsection{Contact models}

We normally keep the contact process at the reference value ($c_i(t) \equiv 1$). To examine how temporal contact variation interacts with the punctuatedness of $b_i(\tau)$, we introduce two contact models that serve as simple test beds: one deterministic and periodic, and one stochastic with heavy-tailed levels.

For the periodic contact model, we assume
%
\begin{equation}
    c_i(t) = 1 - A \cos(\omega t) \quad \forall i
    \label{eq:periodic_c}
\end{equation}
%
with amplitude $A \in [0,1]$ and angular frequency $\omega > 0$. We take $\omega = 2\pi\,\text{day}^{-1}$ as the default, modeling a contact process with a period of one day that reflects higher numbers of contacts during the day and lower numbers at night. Setting $A = 0$ recovers the reference contact process. Here, there is no variation between individuals' contact processes ($C(t) = c_i(t)$), but there is temporal variation within each person's contacts.

For the Lomax contact model, we assume each person's contact modulator is a piecewise-constant process that switches to a new level at the points of a Poisson process with rate $\lambda$, \ie, the sojourn at each level is exponentially distributed with mean $1/\lambda$. The levels are independently drawn from a Lomax (Pareto Type II) distribution with shape $\alpha_P$ and scale $\sigma$,
%
\begin{equation}
    f_{c^*}(c) = \frac{\alpha_P}{\sigma}\Bigl(1 + \frac{c}{\sigma}\Bigr)^{-(\alpha_P + 1)}, \qquad c \geq 0
    \label{eq:lomax_pdf}
\end{equation}
%
with $\alpha_P > 1$. The Lomax distribution has support $[0,\infty)$ and mean $\sigma / (\alpha_P - 1)$, so requiring a mean of 1 pins down
%
\begin{equation}
    \sigma = \alpha_P - 1
    \label{eq:lomax_scale}
\end{equation}
%
leaving a one-parameter family indexed by the tail exponent $\alpha_P$. Under this constraint, the variance of $c^*$ is $\alpha_P / (\alpha_P - 2)$, which is finite for $\alpha_P > 2$ and diverges as $\alpha_P \searrow 2$; smaller $\alpha_P$ corresponds to a heavier upper tail and rarer but more extreme contact levels, while larger $\alpha_P$ makes $c^*$ progressively tighter around 1. In contrast to the periodic model, here each person's contact trajectory is unique, but the population average is constant: $C(t) = E_i[c_i(t)] = 1$.

% ////////////////////////////////////////////////////////////////////////////////////////////

\subsection{Simulation approach}

We investigate the impact of punctuatedness on epidemic dynamics using both closed-form analytic results (derived in the {\bf Supplementary Methods}) and stochastic simulations. The deterministic baseline is obtained by numerically solving the renewal equation with the appropriate $A(\tau)$.

Stochastic simulations use a branching-process-like algorithm in a finite population of $N = 10{,}000$ susceptible individuals with a single initial infectious individual. For each infected person $i$:
%
\begin{enumerate}
    \item A biological infectiousness profile $b_i(\tau)$ is drawn from the time-shifted Gamma model (Eq.~\ref{eq:gamma_timeshift}), and a contact trajectory $c_i(t)$ is drawn from the relevant contact model.
    \item The number of secondary infection attempts is drawn as $X_i \sim \text{Poisson}(R_0)$.
    \item The timing of each attempt relative to infection is drawn from $b_i(\tau)$.
    \item When the contact process is non-uniform ($c_i \neq 1$), Poisson thinning is applied: each attempt at calendar time $t_i + \tau$ is accepted with probability proportional to $c_i(t_i + \tau)$, ensuring that the time-dependent infection rate matches $a_i(t_i, \tau)$.
    \item Each accepted attempt is paired with a uniformly chosen individual from the population: if the recipient is susceptible, infection occurs; otherwise, the attempt fails.
\end{enumerate}
%
Full details are provided in the {\bf Supplementary Methods}. For illustration, we focus on three sets of epidemiological parameters roughly pertaining to influenza-like, SARS-CoV-2 omicron-like, and measles-like pathogens ({\bf Table XX}).
